\documentclass[11pt]{article}

\usepackage[margin=1in]{geometry}
\usepackage{amsmath,amssymb}
\usepackage[hidelinks]{hyperref}
\hypersetup{
  pdftitle={Near-Square-Root NP-Hardness for Euclidean CVP Would Collapse the Exponential Hierarchy},
  pdfauthor={somewhereafter},
  pdfsubject={A conditional exponential-hierarchy barrier for near-square-root Euclidean GapCVP hardness}
}

\newtheorem{theorem}{Theorem}[section]
\newtheorem{corollary}[theorem]{Corollary}
\newtheorem{proposition}[theorem]{Proposition}
\newtheorem{remark}[theorem]{Remark}
\newtheorem{definition}[theorem]{Definition}
\newenvironment{proof}{\par\noindent\textit{Proof.}\ }{\hfill$\square$\par}

\newcommand{\GapCVP}{\mathsf{GapCVP}}
\newcommand{\coAM}{\mathsf{coAM}}
\newcommand{\AM}{\mathsf{AM}}
\newcommand{\qpoly}{\mathsf{qpoly}}
\newcommand{\EXP}{\mathsf{EXP}}
\newcommand{\EH}{\mathsf{EH}}
\newcommand{\PH}{\mathsf{PH}}
\newcommand{\poly}{\operatorname{poly}}
\newcommand{\TAUT}{\mathsf{TAUT}}

\title{Near-Square-Root NP-Hardness for Euclidean CVP\\
Would Collapse the Exponential Hierarchy}
\author{somewhereafter}
\date{August 14, 2026}

\begin{document}
\maketitle

\begin{abstract}
Aggarwal, Bennett, Brakerski, Golovnev, Kumar, Li, Peters,
Stephens-Davidowitz, and Vaikuntanathan proved a quantitative public-coin
protocol placing rank-$D$ Euclidean $\GapCVP_\gamma$ in $\coAM$ with
complexity $2^{O(D/\gamma^2)}$ up to polynomial factors.  We record a
consequence at the near-square-root approximation frontier.  If, for any
fixed $C>0$, Euclidean
\[
  \GapCVP_{\sqrt D/\log^C D}
\]
is NP-hard under polynomial-time promise-preserving many-one reductions,
then the exponential hierarchy collapses to its second level.  The proof
combines the quantitative lattice protocol with the resource-scaled
Boppana--H\aa stad--Zachos theorem of Pavan, Selman, Sengupta, and
Vinodchandran.

More generally, for a reduction from inputs of length $n$ producing rank
$D(n)$ and approximation $\gamma(D(n))$, the same conclusion follows whenever
$D(n)/\gamma(D(n))^2=\log^{O(1)}n$.  Writing
$\gamma(D)=\sqrt D/f(D)$ exposes the exact protocol frontier: every fixed
polylogarithmic loss $f(D)=\log^{O(1)}D$ is conditionally excluded.  For
rank-regular reductions, the first dimension-only regime not excluded by this
argument is the superpolylogarithmic-loss strip
$f(D)\notin\log^{O(1)}D$.  The genuinely near-square-root portion of that
region additionally has $f(D)=D^{o(1)}$.  This is a barrier theorem, not an
unconditional lower bound and not a construction in the surviving strip.
\end{abstract}

\section{Introduction}

The Euclidean closest vector problem asks for the distance from a target
$t\in\mathbb{Q}^D$ to a full-rank lattice $L\subseteq\mathbb{R}^D$.
Its promise version $\GapCVP_\gamma$ distinguishes
\[
  \operatorname{dist}(t,L)\le r
  \qquad\text{from}\qquad
  \operatorname{dist}(t,L)>\gamma(D)r.
\]
The square-root scale is a persistent boundary in the complexity of lattice
problems, beginning with the classical coAM barrier of Goldreich and
Goldwasser~\cite{GG00}.  A particularly sharp prospective hardness factor is
\[
  \gamma_C(D)=\frac{\sqrt D}{\log^C D},                 \tag{1}
\]
where $C>0$ is fixed.  This factor is $D^{1/2-o(1)}$ and is asymptotically
larger than $D^c$ for every fixed $c<1/2$.
It should be compared with Mira's formalized NP-hardness theorem at $D^c$ for
every fixed $c<1/2$~\cite{Mira26}.

Our main observation is that (1) is already inside a universal
quasipolynomial-time $\coAM$ regime.  This fact has a stronger consequence
than a construction-specific obstruction: polynomial-time NP-hardness at
(1) would collapse the exponential hierarchy.  No classification of lattice
gadgets, codes, Fourier witnesses, or source compilers is needed.

The ingredients are published.  Aggarwal et al.~\cite{ABBGKLPV26} prove the
time--approximation tradeoff
\[
  \GapCVP_\gamma\in\coAM\text{-TIME}
  \left(2^{O(D/\gamma^2)}\poly(M)\right),               \tag{2}
\]
where $M$ is the target-instance encoding length.  Pavan et
al.~\cite{PSSV07} prove a resource-scaled extension of the
Boppana--H\aa stad--Zachos collapse theorem~\cite{BHZ87}: quasipolynomial
Arthur--Merlin protocols for coNP collapse the exponential hierarchy to its
second level.  Substituting (1) into (2) gives complexity
$2^{O(\log^{2C}D)}\poly(M)$, which is quasipolynomial under any polynomial-time
reduction.

The contribution of this note is threefold.
\begin{enumerate}
  \item We state and prove the resulting architecture-independent conditional
  impossibility theorem for near-square-root $\GapCVP$ hardness.
  \item We give the source-sensitive form
  $D(n)/\gamma(D(n))^2=\log^{O(1)}n$, avoiding an implicit assumption that
  target rank and source length are polynomially equivalent.
  \item We isolate the first asymptotic strip not covered by the
  quasipolynomial-protocol argument and separate it from the stronger AMETH
  barrier already implicit in~\cite{ABBGKLPV26}.
\end{enumerate}

The theorem is conditional.  It does not show unconditionally that the
reduction cannot exist.  Rather, a correct reduction would itself establish
a major and unexpected collapse of complexity classes.

\section{Preliminaries}

\subsection{Promise hardness}

We use deterministic polynomial-time promise-preserving many-one reductions.
If $A$ is a language and $\Pi=(\Pi_Y,\Pi_N)$ is a promise problem, a reduction
$R$ satisfies
\[
  x\in A\Rightarrow R(x)\in\Pi_Y,
  \qquad
  x\notin A\Rightarrow R(x)\in\Pi_N.
\]
We call $\Pi$ NP-hard if such a reduction exists from an NP-complete language.

Let $n$ denote source length, $M(n)$ the bit length of the emitted lattice
instance, and $D(n)$ its rank.  Polynomial time implies
\[
  M(n)\le n^{O(1)},\qquad D(n)\le M(n).                 \tag{3}
\]
We do not assume a polynomial lower bound on $D(n)$ unless stated.

\subsection{Quasipolynomial and exponential resources}

Write
\[
  \qpoly(n)=2^{(\log n)^{O(1)}}.
\]
The class $\AM_{\qpoly}$ consists of languages with constant-round public-coin
protocols whose verifier time, randomness, and total communication are
quasipolynomial.  Define $\PH_{\qpoly}$ by allowing a fixed number of
alternations with quasipolynomial resources.

The exponential hierarchy is
\[
  \EH=\bigcup_{k\ge1}\Sigma_k^{\EXP}.
\]
We use the following published collapse theorem.

\begin{theorem}[Pavan--Selman--Sengupta--Vinodchandran
\cite{PSSV07}]                                                    \label{thm:pssv}
If $\mathsf{coNP}\subseteq\AM_{\qpoly}$, then
\[
  \PH_{\qpoly}\subseteq\AM_{\qpoly}
\]
and, by padding,
\[
  \EH\subseteq\AM_{\EXP}.
\]
The standard containment
$\AM_{\EXP}\subseteq\Sigma_2^{\EXP}\cap\Pi_2^{\EXP}$, together with
$\Sigma_2^{\EXP}\subseteq\EH$, therefore gives
\[
  \EH=\AM_{\EXP}\subseteq\Pi_2^{\EXP}.
\]
Consequently the exponential hierarchy collapses to its second level.
\end{theorem}

This strengthens an earlier theorem of Selman and Sengupta that obtained only
a third-level collapse.  The distinction is important: the later argument
extends the Boppana--H\aa stad--Zachos proof to quasipolynomial resource bounds
before applying padding.

\section{The quantitative \texorpdfstring{$\coAM$}{coAM} protocol}

Aggarwal et al. prove the following tradeoff.

\begin{theorem}[Aggarwal et al.]                                  \label{thm:coam}
Corollary 3.3 of~\cite{ABBGKLPV26} states that for every
$\gamma=\gamma(D)\ge1$, rank-$D$ Euclidean
$\GapCVP_\gamma$ has a constant-round $\coAM$ protocol of complexity
\[
  2^{O(D/\gamma^2)}\poly(M),                              \tag{4}
\]
where $M$ is the input encoding length.
\end{theorem}

The underlying private-coin protocol performs
\[
  N_2=10D^{3/2}(1-1/\gamma^2)^{-(D+1)/2}                 \tag{5}
\]
independent trials.  In each trial Arthur chooses a hidden uniform bit and a
uniform point in one of two translated Euclidean balls; Merlin must identify
the hidden translate.  Generic private-to-public-coin and round-reduction
transformations incur polynomial overhead, preserving the exponent in (4).

For $\gamma(D)=\sqrt D/f(D)$, (4) becomes
\[
  2^{O(f(D)^2)}\poly(M).                                  \tag{6}
\]
Equation (6) is the only lattice-specific input needed below.

\section{The hierarchy-collapse barrier}

We first state the source-sensitive theorem.

\begin{theorem}[General reduction barrier]                         \label{thm:general}
Suppose there is a deterministic polynomial-time promise-preserving many-one
reduction from $\mathsf{SAT}$ instances of length $n$ to Euclidean
$\GapCVP_{\gamma}$ instances of rank $D(n)$ and encoding length $M(n)$.
If, for some fixed $a$,
\[
  \frac{D(n)}{\gamma(D(n))^2}\le (\log n)^a              \tag{7}
\]
for all sufficiently large $n$, then
\[
  \EH=\AM_{\EXP}\subseteq\Pi_2^{\EXP}.
\]
In particular, the exponential hierarchy collapses to its second level.
\end{theorem}

\begin{proof}
By Theorem~\ref{thm:coam}, the target promise problem has a $\coAM$ protocol
of complexity
\[
  2^{O(D(n)/\gamma(D(n))^2)}\poly(M(n)).
\]
Equations (3) and (7) make this quantity quasipolynomial in $n$.  Composing
the verifier with the deterministic reduction gives
\[
  \mathsf{SAT}\in\coAM_{\qpoly},
\]
and complementing gives $\TAUT\in\AM_{\qpoly}$.  Since $\TAUT$ is
coNP-complete under polynomial-time many-one reductions,
\[
  \mathsf{coNP}\subseteq\AM_{\qpoly}.
\]
The conclusion follows from Theorem~\ref{thm:pssv}.
\end{proof}

\begin{corollary}[Fixed-polylogarithmic loss]                       \label{cor:polylog}
For every fixed $C>0$, if
\[
  \GapCVP_{\sqrt D/\log^C D}
\]
is NP-hard under polynomial-time promise-preserving many-one reductions, then
the exponential hierarchy collapses to its second level.
\end{corollary}

\begin{proof}
Here
\[
  \frac{D}{\gamma(D)^2}=\log^{2C}D.
\]
Since $D(n)\le M(n)\le n^{O(1)}$, this is $\log^{O(1)}n$, so
Theorem~\ref{thm:general} applies.
\end{proof}

\begin{corollary}[Conditional impossibility]
Assume that the exponential hierarchy does not collapse to its second level.
Then, for every fixed $C>0$, Euclidean
$\GapCVP_{\sqrt D/\log^C D}$ is not NP-hard under deterministic
polynomial-time promise-preserving many-one reductions.
\end{corollary}

This conclusion is independent of the internal form of the reduction.  It
applies equally to reductions built from codes, PCPs, Fourier statistics,
group actions, geometric gadgets, or any other polynomial-time compiler.

\section{The exact asymptotic frontier of this argument}

Write the approximation factor as
\[
  \gamma(D)=\frac{\sqrt D}{f(D)},                         \tag{8}
\]
where $f(D)$ is the loss from the square-root scale.  Theorem~\ref{thm:general}
becomes the following exact criterion.

\begin{proposition}[Protocol-frontier criterion]                   \label{prop:frontier}
A polynomial-time reduction with loss $f$ implies the second-level
exponential-hierarchy collapse whenever
\[
  f(D(n))^2=\log^{O(1)}n.                                 \tag{9}
\]
In particular, every dimension-only bound
\[
  f(D)=\log^{O(1)}D                                      \tag{10}
\]
satisfies (9) for every polynomial-time reduction.
\end{proposition}

The converse is deliberately not claimed.  If $f(D)$ is superpolylogarithmic
in $D$, condition (9) may still hold when the reduction emits unusually small
rank relative to source length.  To state a clean dimension-only boundary,
one may impose the natural rank-regularity condition
\[
  D(n)\ge n^{\delta}                                     \tag{11}
\]
for some fixed $\delta>0$.

\begin{corollary}[Rank-regular boundary]
Under (11), the protocol in (4) has quasipolynomial complexity precisely when
\[
  f(D)=\log^{O(1)}D
\]
along the output ranks, up to the usual eventual-bound interpretation.
Thus the first dimension-only region not covered by the EH-noncollapse
argument consists of losses not bounded by any fixed polylogarithm.  The
near-square-root portion of this region is
\[
  \log^C D=o(f(D))\text{ for every fixed }C,
  \qquad f(D)=D^{o(1)},
  \qquad
  \gamma(D)=\frac{\sqrt D}{\log^{\omega(1)}D}.           \tag{12}
\]
\end{corollary}

There is no smallest superpolylogarithmic function.  Equation (12) describes
a boundary class, not one canonical approximation factor.  Examples include
\[
  f(D)=\exp((\log\log D)^{1+\varepsilon})
  \quad\text{and}\quad
  f(D)=\exp(\sqrt{\log D}).
\]
Both give $\gamma(D)=D^{1/2-o(1)}$, which is asymptotically larger than
$D^c$ for every fixed $c<1/2$, while the protocol in (6) is
superquasipolynomial.

Failure of the present argument in (12) is not evidence that hardness exists
there.  It says only that the published quantitative $\coAM$ protocol no
longer places the composed source problem in $\AM_{\qpoly}$.

\section{The stronger AMETH boundary}

Aggarwal et al. define AMETH as the assertion that, for some $\varepsilon>0$,
$3$-$\TAUT$ has no constant-round public-coin protocol of complexity
$2^{\varepsilon n}$.  The same composition gives a stronger but less
structural conditional barrier.  Here, unlike in Sections 2--5, $n$ denotes
the number of variables.  We restrict the proposed reduction to $3$-CNF
formulas with $O(n)$ clauses and let $D(n)$ be the largest emitted rank on
those inputs.  Their encoding length is $O(n)$; the standard sparsification
lemma~\cite{IPZ01} lifts any $2^{o(n)}$ protocol for this sparse restriction to a
$2^{\delta n+o(n)}$ protocol for general $3$-$\TAUT$, for every fixed
$\delta>0$.

\begin{proposition}[AMETH barrier]                                \label{prop:ameth}
Under AMETH, no reduction as in Theorem~\ref{thm:general} can satisfy
\[
  \frac{D(n)}{\gamma(D(n))^2}=o(n).                       \tag{13}
\]
on all $3$-CNF inputs with $n$ variables and $O(n)$ clauses.
Equivalently, for (8), it cannot satisfy
\[
  f(D(n))^2=o(n).                                         \tag{14}
\]
\end{proposition}

\begin{proof}
Equations (4) and (13) give a $2^{o(n)}$ constant-round protocol for sparse
$3$-$\TAUT$.  Apply the sparsification lemma and run the resulting sparse protocols in
parallel.  For every fixed $\delta>0$, this gives a constant-round protocol
of complexity $2^{\delta n+o(n)}$ for general $3$-$\TAUT$.  Taking
$\delta$ below the constant in AMETH yields a contradiction.
\end{proof}

If $D(n)\le n^K$ for fixed $K$ and $f(D)=D^{o(1)}$, then
$f(D(n))^2=n^{o(1)}=o(n)$.  Hence AMETH excludes the entire
$D^{1/2-o(1)}$ approximation regime for polynomial-output reductions, not
only fixed-polylogarithmic losses.  The superpolylogarithmic strip (12) is
therefore open only relative to the weaker EH-noncollapse assumption; it is
already forbidden by AMETH.

For a fixed-power factor $\gamma(D)=D^c$, the protocol exponent is
$D^{1-2c}$.  If $D(n)\le n^K$, AMETH forces the efficiency ceiling
\[
  K(1-2c)\not<1,
  \qquad\text{or heuristically}\qquad
  c\le\frac12-\frac{1}{2K},                              \tag{15}
\]
unless constants or lower-order behavior keep (13) from holding.  Thus
hardness for every fixed $c<1/2$ can coexist with the barrier when approaching
$1/2$ requires an unbounded output-degree parameter; a uniform
$D^{1/2-o(1)}$ theorem cannot.

\section{Why the protocol does not give a deterministic algorithm}

The conclusion of Corollary~\ref{cor:polylog} is an interactive-proof
collapse, not a quasipolynomial-time algorithm for $\GapCVP$.  This distinction
cannot be removed by naively enumerating Arthur's random coins.

The private protocol underlying (4) uses $N_2$ independent hidden bits even
before accounting for the random convex-body samples.  Therefore its written
implementation consumes at least $N_2$ random bits.  At the near-square-root
endpoint, substituting $\gamma=\sqrt D/f(D)$ into (5) gives
\[
  \log N_2
  =\log 10+\frac32\log D+\frac12f(D)^2
   +O\!\left(\frac{f(D)^2+f(D)^4}{D}\right).             \tag{16}
\]
Thus $N_2=2^{\log^{O(1)}D}$ for polylogarithmic $f$, and exhaustive
enumeration costs
\[
  2^{\Omega(N_2)},
\]
not quasipolynomial time.  Corollary 3.3 of~\cite{ABBGKLPV26} invokes generic
Goldwasser--Sipser~\cite{GS89} and Babai--Moran~\cite{BM88}
transformations with polynomial overhead;
it does not replace the experiment by a pairwise-independent test with a
logarithmic seed.

More precisely, let $\ell$ be the amplified private random-tape length in the
Goldwasser--Sipser simulation.  Its explicit public-coin implementation uses
full random linear maps and at least $2\ell^2$ public random bits at an
intermediate stage.  Since $\ell\ge N_2$, enumerating those public strings
already costs $2^{\Omega(N_2^2)}$.  The final Babai--Moran transformation is
stated only with polynomial resource overhead and supplies no alternative
polylogarithmic seed bound.  Pairwise-universal hashing in the proof therefore
does not amount to a short-seed derandomization.

Nor can a standard pseudorandom generator simply be substituted.  In a
public-coin protocol Merlin's response depends on the revealed coins, and
soundness quantifies over all prover strategies.  Fooling a fixed
space-bounded computation is not by itself a derandomization of this
interactive maximum.  Accordingly, this note claims
$\mathsf{coNP}\subseteq\AM_{\qpoly}$ under the hypothetical hardness
reduction, but not $\mathsf{coNP}\subseteq\mathsf{QP}$ and not a
deterministic quasipolynomial-time algorithm for $\GapCVP$.

One may obtain the standard nonuniform containment
\[
  \coAM_{\qpoly}\subseteq\mathsf{coNP}/\qpoly
\]
by fixing a quasipolynomial-size advice list of public strings after
amplification and a union bound over all inputs of a given length.  The advice
slash is essential; this is not a uniform seed-enumeration algorithm and is
not used in our main theorem.

\section{Relation to fixed-exponent hardness}

Mira (\texttt{@\_Mira\_\_\_Mira\_}) proves NP-hardness at $D^c$ for every
fixed $c<1/2$~\cite{Mira26}.  This theorem does not diagonalize to (1): each
fixed exponent corresponds to polynomial loss
\[
  f(D)=D^{1/2-c},
\]
whereas (1) has only polylogarithmic loss.  A family of separate reductions,
one for each fixed $c$, need not yield a single polynomial-time reduction with
$c=c(D)\to1/2$.  Typically the polynomial degree of the output size depends on
$c$ and becomes unbounded as $c\to1/2$.

Under EH noncollapse, Corollary~\ref{cor:polylog} says that this distinction is
not merely a defect of one known construction: no polynomial-time Karp
reduction can reach the fixed-polylogarithmic-loss endpoint.  It does not rule
out an improvement into the superpolylogarithmic strip (12).

\section{Conclusion}

The quantitative $\coAM$ protocol for Euclidean $\GapCVP$ and the
resource-scaled interactive-proof collapse theorem together give a universal
conditional boundary:
\[
  \boxed{
  \text{NP-hardness at }\sqrt D/\log^{O(1)}D
  \Longrightarrow
  \text{the exponential hierarchy collapses to level two}.}
\]
This makes any such hardness proof simultaneously a hierarchy-collapse proof.
Assuming EH noncollapse, the fixed-polylogarithmic-loss endpoint is closed.

The next unresolved asymptotic target, for rank-regular reductions and under
that assumption alone, is hardness with a loss that is simultaneously
superpolylogarithmic and subpolynomial, schematically
\[
  \frac{\sqrt D}{\log^{\omega(1)}D}.
\]
This strip lies strictly beyond every fixed power $D^c$, $c<1/2$, but outside
the quasipolynomial-$\coAM$ zone.  Establishing hardness there would require a
new reduction; excluding it would require an assumption or upper bound
stronger than EH noncollapse.

\bibliographystyle{alpha}
\bibliography{references}

\end{document}
